\section*{結言}
今回の中心は、フロンティアモデルの能力向上と、その利用範囲を制御する動きが同時に進んだことにある。Gemini 4 Carbonはサイバーセキュリティと長期タスクの性能を掲げつつ、初期アクセスを防御側に絞る。Genie 4を巡る議論も含め、モデルの新機能だけでなく、誰がどの条件で使えるかが重要な評価軸になっている。

公開モデルでは、Cloverの拡散生成とLFM3の小型化が異なる効率化の方向を示した。Articul8のBiviaやVangrid関連の投稿は、AIを業務判断や現実世界のデータ収集へ組み込む関心を映す。ただし、発表者が示す速度や価格、実用性は利用条件に左右されるため、独立した評価が求められる。

一方、Amazon Novaの再編報道やDeepSeek V5のリークには一次確認の不足が残る。観測窓を越える投稿も一部含まれており、速報で見られた話題と、当該24時間に確認された事実は区別する必要がある。次の焦点はCarbonの提供拡大、公開モデルの実測比較、未確認情報に対する公式の説明となる。
